% 全体ゼミ資料：物理モデルの自動構築のための数式集合の検索 — 変数を記号でなく説明文で照合する試み
% ビルド: このディレクトリで latexmk（latexmkrc: platex → pbibtex → dvipdfmx）。外部画像なし。
% 文献: progress_report_2026-10-08_var_description.bib（natbib + unsrtnat。本文は \citet / \citep）。
% 改訂履歴:
%   2026-10-08 初版（本人指示「このデザインを簡単な進捗報告としてまとめて。前のやり方と今回のやり方の比較。2 ページに収まって」）。設計だけの版。
%   2026-10-08 別セッション（論文用 Methods）が本人の指示で Abstract 冒頭の目的の文を 2 文に書き換えた（修論とそろえた文。書き直しても残す）。
%   2026-10-09 本人指示「全ての生成・実験が終わってから書こう」に従い、本実験と追加実験の結果を入れて書き直した。
%   2026-10-12 本人指示「作文技術チェックリストを踏まえて改善。研究室の他のグループ（他の分野）向けなので、この研究に限る難しいことは程よく略して簡単に」:
%     - τ の記号と決め方の名前、1 対 1 の対応、第 1 段の重み、設定 A・B の呼び名、写した説明の参考の行を本文から外し、埋め込みは言葉で説明した。
%     - Abstract の比較を 8/31 報告の改訂（最良の重みの baseline、0.607）にそろえ、第 1 段の重み (0.3, 0.7) の提案手法（0.780）と比べた。
%     - 正解式を定義し、有意の横に検定名と p 値を書き、0 でない値を 0.000 と書かない（0.0002）。
%     2 ページに収めるため Appendix は置かない。
%   2026-10-12 本人指示: ヘッダの日付を 10 月 12 日に、確認者の行を削除（Abstract の長さは問題ない）。
%   2026-10-12 本人指示: Abstract を本人の文にした（目的の 2 文目「提示した数式からモデルを組み立てて確かめるのは利用者である。」を削除）。
%   2026-10-12 本人指示: 2.1 の「実験設定の妥当性の確認と、それに合わせた実験のやり直し」を削除。議論したい点 1（ProcessBERT）は「やらないと分からないので意見を求めても意味がない」ため、今後の計画（試すこと）に移した。
%   2026-10-12 本人指示「2.2 未回答だった質問への対応 節。なしなら消して」: 2.2 節を見出しごと削除（2.1 節だけを残す。9/27 版で本人が消したのと同じ形）。
%   2026-10-13 本人指示: 2.1 の前置き「前回のゼミ（9 月 29 日）で挙げた計画の状況は次のとおりである。」と、補足「（3 節）」（1 節と 2.1 の 2 か所）を削除。
%   2026-10-13 本人指示「3.1 何が問題かを考え直して。きっかけは文脈で精度を上げると主張したのに文脈の特徴量が効かなかったから、文脈が効く場面を作りたい」: 3.1 を、文脈の特徴量が効かなかったこと→評価用ケースが文献と同じ記号を使うためと考えられる→記号が違う場面では文脈が要る→記号が違うケースを作り変数を文脈で照合、の順に書き直した。
%   2026-10-13 本人指示「しきい値の具体値とその値にした理由も書いて」: 3.2 にしきい値（約 0.89）と決め方（同じ変数と別の変数の説明の組を最もよく分ける値 約 0.85 から、調整用ケースの Recall@K で上げ幅を選んだ）を追記。3.3 の該当文を「上げ幅」に合わせた。
%   2026-10-13 本人指示「0.85 から上下ずらして一番精度の高い値を選んだと書けば良い」: 刻みと上げ幅の記述を外し「ずらして試し、最も高い値（約 0.89）を選んだ」に。
%     ただし選び直しの格子は上げる方向（+0〜+0.08）だけなので「上下」とは書かない（下げた値は seed 42 のテスト用ケースで −0.02 を試しただけ）。
%   2026-10-13 本人指示「いらない言葉が多い」: 3.3 冒頭を「評価には 1,823 ケースを使った。各ケースには、モデルに必要な式（正解式）を決めてある。」に。
%   2026-10-13 本人指示「Table 2 の右の列は A で直して」: どのケースも組み合わせて作ったので「組み合わせて作ったケース」をやめ、列名を「正解式が多いケース」、本文とキャプションに作り方（微分方程式から規則で式を足す）と正解式の平均本数を書いた。
%     平均本数: training_cases.json の評価 1,823 件で正解式 7,106 本（平均 3.90）、合成 1,000 件で 5,500 本（平均 5.50）。
%   2026-10-13 本人指示「両方入れて」: 3.2 の利用者の説明文の作り方を分かりやすく書き換え、3.3 にケースの出所（文献のモデルそのままではなく、DB の式を LLM や規則で組み合わせた）を追記。
%     出所: training_cases.json と生成スクリプト。LLM が組を選んだ 323 件（単一出典 92 Gemini 2.5 Pro、出典横断 231 Claude）、プログラム・規則 1,500 件（出典横断 500、合成 1,000）。
%     2 ページに収めるため、3.1 の最後の文、3.2 の利用者の説明文の文、3.3 の選び直しの括弧を短くした。
% 数値の出典:
%   - しきい値: experiments/var_desc_A.json の tau.D1（Youden、seed ごとに 0.8485–0.8503）、var_desc_B.json は 0.8433–0.8485。
%     選んだ値は var_desc_tuned_{A,B}.json の selected.D1: Δτ は A・B とも 10 seed 中 9 で +0.04、1 で +0.02。τ の平均 A 0.888（0.869–0.890）・B 0.885。
%     正例・負例の組の定義は var_conditions.calibration_pairs（正解式の中の同じ記号／同じ正解式の中の別の記号）。
%   - 0.607 / 0.780: docs/progress_report_2026-08-31_methods_results.tex の表（設定 A、baseline の最良の重み (0.4, 0.6) と提案手法 (0.3, 0.7)）。
%     0.780 は experiments/var_desc_A.json の S0（記号による照合、10 seed の平均 0.7799）と同じ。
%   - Table 1 の類似度: intfloat/e5-base-v2、両側に接頭辞 "query: "、平均プーリング＋L2 正規化（本実験と同じ）で計算した。
%     ケース側の説明は ms_cstr_heat_001_rv4 の LLM の言い換え、DB 側は正解式 handbook_chemeng__eq_3 と handbook_heat_transfer__eq_5 の説明、
%     頻度因子は DB の記号 A の説明 "Pre-exponential factor in Arrhenius' equation"。
%   - Table 2: experiments/var_desc_{A,B}.json（S1、S1-r50、S1-r100、D1）と var_desc_tuned_{A,B}.json（D1-tuned）の per_seed（10 seed の平均 ± 標本標準偏差）。
%     本文の「説明文で照合」は、しきい値を調整用（検証用）ケースで選んだ D1-tuned。最初のしきい値（0.85 前後）の D1 は 0.326。
%   - 差と検定: experiments/var_desc_tuned_stats.json（D1-tuned − S1 は −0.338、D1-tuned − S1-r50 は +0.227、Wilcoxon p = 0.00195、t 検定 p < 0.001）。
%   - 第 1 段の候補に正解式が残る割合 0.969 / 0.710: experiments/var_desc_coverage.json（全ケース、文章 0.3・変数 0.7、最初のしきい値）。
\documentclass[a4j,dvipdfmx,10pt]{jresume}
\usepackage[T1]{fontenc} % タイプライターの一部文字（_等）が数式フォントになることを避ける

\usepackage{graphicx}
\usepackage[top=25truemm,bottom=25truemm,left=25truemm,right=25truemm]{geometry}
\usepackage{amsmath, amssymb, array, bm}
\usepackage{booktabs}
\usepackage{url}
\usepackage{xcolor} % コメント用（blue!55!black などの色の指定にも使う）
\usepackage{ulem} % コメント用
\normalem
\usepackage[numbers,square,sort&compress]{natbib}
\usepackage{hyperref}
\usepackage{doi}
% 本報告で追加したパッケージ（箇条書きの間隔）
\usepackage{enumitem}

\setlength{\abovecaptionskip}{0pt}
\setlength{\belowcaptionskip}{0pt}
\setlength{\textfloatsep}{8pt plus 2pt minus 2pt} % 2 ページに収めるため、表と本文の間を詰める
\setlength{\floatsep}{6pt plus 2pt minus 2pt}
\hypersetup{colorlinks=true, linkcolor=blue!55!black, urlcolor=blue!55!black, citecolor=blue!55!black}

%%%%%%%%%%%%%%%%%%%%%%%%%%%%%%%%%%%%%
%%%%%%  User-defined commands  %%%%%%
%%%%%%%%%%%%%%%%%%%%%%%%%%%%%%%%%%%%%
\renewcommand{\refname}{References}
\renewcommand{\abstractname}{Abstract}

%% コメント用
\newcommand{\kato}[1]{\textcolor[rgb]{0.9,0,0}{#1}}
\DeclareRobustCommand{\katoerase}{\bgroup\markoverwith{\textcolor[rgb]{0.9,0,0}{\rule[.5ex]{2pt}{0.8pt}}}\ULon}

\newcommand{\rk}{\textrm{Recall@}K}
\setlist{nosep}
\renewcommand{\bibfont}{\footnotesize}
\setlength{\bibsep}{1pt}
%%%%%%%%%%%%%%%%%%%%%%%%%%%%%%%%%%%%%

\pagestyle{myresume}

\begin{document}
%%%%%     DATE     %%%%%
\西暦 % この行の有無で年表示が変わる
\markright{{\footnotesize 研究進捗報告 \hfill 2026 年 10 月 12 日}}

%%%%%     TITLE     %%%%%
\medskip
\begin{center}
  {\Large 物理モデルの自動構築のための数式集合の検索\\ 変数を記号でなく説明文で照合する試み}
\end{center}

%%%%%     AUTHOR     %%%%%
\begin{flushright}
  報告者：M2\; 一洋 \\
\end{flushright}

\begin{abstract}
本研究の目的は、複数の情報源から収集した数式の中から、作りたいモデルに使える可能性の高い数式を提示する手法を開発することである。作成するモデルは、説明文と入出力変数で表される。提案手法は 2 段階の処理から構成される。第 1 段では、数式を説明文と式周辺の文章（\textbf{文脈}）の近さと入出力変数の重複率で採点し、上位 $N_1$ 件に絞る。第 2 段では、10 個の特徴量を持つ多層パーセプトロンで並べ直す。提案手法を 11{,}146 式と評価用の 1{,}823 ケースで評価した結果、上位 $K$ 件に入った正解式の割合 $\rk$（$K$ は正解式の数）は、第 1 段と同じ点数で全数式を並べる古典的な情報検索手法（最良の重み）の 0.607 から 0.780 に上がった。並べ直しでは、主に変数の重複率に基づく特徴量が効き、文脈の近さはほとんど効かなかった。
\end{abstract}

%%%%%%     BODY START     %%%%%
\section{進捗と議論のポイント}
\begin{itemize}[leftmargin=1.5em]
  \item 変数を記号（例：$C_{A0}$）ではなく、意味を表す説明文（例：「原料中の成分 A の濃度」）で照合する方法を試した。利用者が文献と同じ記号を使う場合は、記号で照合する従来の方法の方が、正解式を多く上位に並べた（$\rk$ で 0.780 対 0.442）。
  \item 利用者が文献と違う記号を使う場合を模擬すると、記号で照合する方法は大きく下がった（半分の記号を変えると 0.215）。説明文で照合する方法は記号を使わないので影響を受けず、これを上回った。
  \item \textbf{議論したい点}：記号が文献と違っても正解式を探せることを、変数の説明（文脈）を扱える本研究の利点として主張してよいか、意見をいただきたい。
\end{itemize}


\section{前回の振り返り}\label{sec:review}
\subsection{計画の達成状況}
\begin{itemize}[leftmargin=1.5em]
  \item \textbf{文脈が効く場面を示す評価の設計}：達成。記号を変えたケースで 2 つの照合を比べた。
\end{itemize}


\section{変数を記号でなく説明文で照合する}\label{sec:method}
\subsection{何が問題か}
本研究は、式周辺の文章（文脈）を扱えることを先行研究との違いとしてきた。しかしこれまでの評価では、文脈の近さに基づく特徴量はほとんど効かず、正解式を上位に並べたのは主に変数の重複率に基づく特徴量であった。原因の 1 つは評価用のケースの作り方にあると考えられる。評価用のケースは入出力変数を文献と同じ記号で与えるので、記号が一致する変数を数えるだけで正解式を見つけられ、文脈を使う必要がない。しかし実際の利用者は文献の記号を知らないことが多く、同じ量を違う記号で書くことがある（例：原料中の成分 A の濃度を $C_{A0}$ ではなく $c_{A,\mathrm{in}}$ と書く）。また、同じ記号 $A$ が、ある文献では伝熱面積を、別の文献では反応速度定数の頻度因子を指すこともある。このような場面では記号の照合が働かず、変数の意味を文脈から読み取る必要がある。これまで文脈は作りたいモデルの説明文との近さにだけ使い、変数の照合には使っていなかった。そこで、記号が文献と違うケースを作り、変数を文脈で照合した。

\subsection{今回の方法}
変数を記号ではなく、その意味を表す説明文で照合する。2 つの説明文の近さは、汎用の文章埋め込みモデル E5（EmbEddings from bidirEctional Encoder rEpresentations）~\citep{E5}で各文を数値のベクトルに変え、その向きの近さ（コサイン類似度。同じ向きなら 1）で測り、しきい値以上なら同じ変数とみなす。しきい値は約 0.89 とした。まず、正解式の中で同じ記号の変数の説明の組を「同じ変数」、別の記号の変数の説明の組を「別の変数」として、2 つを最もよく分ける値（正しく同じとみなす割合と誤って同じとみなす割合の差が最大になる値）を、テストに使わないケースで求めた（約 0.85）。しかし、変数の組を最もよく分ける値が、検索の精度を最も高くするとは限らない。そこで、この値からしきい値をずらして試し、テストに使わないケースで $\rk$ が最も高くなる値（約 0.89）を選んだ。式の側にはデータベースの各式に付いている変数の説明を使う。評価用のケースには利用者が書いた変数の説明がない。正解式の説明をそのまま使うと式の側と同じ文になり、必ず一致する。利用者は自分の言葉で書くので、正解式の説明を大規模言語モデル（large language model, LLM）で言い換えて使った。2 つの照合の違いの例を Table~\ref{tab:compare} に示す。

\begin{table}[t]\centering
\footnotesize
\caption{Matching by symbols and by descriptions in three typical cases (values: cosine similarities).}
\label{tab:compare}
\begin{tabular}{@{}>{\raggedright\arraybackslash}p{0.30\linewidth}>{\raggedright\arraybackslash}p{0.25\linewidth}>{\raggedright\arraybackslash}p{0.39\linewidth}@{}}
\toprule
例 & 記号で照合 & 説明文で照合 \\
\midrule
1. 同じ量を違う記号で書く（$C_{A0}$ と $c_{A,\mathrm{in}}$） & 一致しない（見逃す） & 記号を使わないので影響を受けない（入口の濃度の説明どうしで 0.929） \\
\addlinespace
2. 同じ記号が別の量を指す（伝熱面積の $A$ と頻度因子の $A$） & 一致する（誤って結ぶ） & 区別できる（伝熱面積どうしは 0.893、伝熱面積と頻度因子は 0.752） \\
\addlinespace
3. 似ているが別の量（入口の濃度と出口の濃度） & 記号が違えば区別できる & 区別できない（入口どうしは 0.929、入口と出口は 0.933） \\
\bottomrule
\end{tabular}
\end{table}

\subsection{実験と結果}\label{sec:result}
評価には 1{,}823 ケースを使った。各ケースには、モデルに必要な式（正解式）を決めてある。ケースは文献のモデルそのものではなく、データベースの式を LLM や規則で組み合わせて作った。ケースを学習用、調整用、テスト用に分ける分け方を 10 通り変え、テスト用ケースの $\rk$ の平均を比べた。利用者が文献と違う記号を書く場面は、テスト用ケースの入出力変数の記号を半分またはすべて、LLM が提案した別の記号に変えて模擬した（学習用ケースとデータベースは元の記号のまま）。説明文で照合するときのしきい値と第 1 段の点数の重みは、調整用ケースで $\rk$ が最大になる値を選んだ（最初の値では 0.326、選び直すと 0.442）。記号で照合する方法の重みも同じ方法で選んだが、もとの値が選ばれた。

\begin{table}[t]\centering
\footnotesize
\caption{Recall@K on the test cases (mean $\pm$ SD, 10 splits). Right: the 1,000 cases built by a rule that starts from a differential equation and adds equations (5.5 correct equations on average; 3.9 for all cases).}
\label{tab:result}
\begin{tabular}{@{}lcc@{}}
\toprule
照合の方法 & 全ケース & 正解式が多いケース \\
\midrule
記号で照合（記号は文献と同じ） & 0.780 $\pm$ 0.046 & 0.745 $\pm$ 0.040 \\
記号で照合（半分の記号を変える） & 0.215 $\pm$ 0.024 & 0.235 $\pm$ 0.028 \\
記号で照合（すべての記号を変える） & 0.0002 $\pm$ 0.0004 & 0.0003 $\pm$ 0.0005 \\
説明文で照合（記号を使わない） & 0.442 $\pm$ 0.046 & 0.371 $\pm$ 0.024 \\
\bottomrule
\end{tabular}
\end{table}

結果を Table~\ref{tab:result} に示す。記号が文献と同じなら、記号で照合する方が正解式を多く上位に並べた（全ケースで 0.780 対 0.442）。一方、半分の記号を変えると、記号で照合する方法は 0.215 に下がり、説明文で照合する方法を下回った（差は全ケースで 0.227）。すべての記号を変えると 0.0002 とほぼ 0 になった。変えた記号がデータベースの別の式でも使われていて、関係のない式を上位に押し上げるためと考えられる。2 つの差はどちらも、10 通りの分け方すべてで同じ向きであった（Wilcoxon の符号順位検定で $p = 0.002$。10 組で取りうる最小の値。対応のある $t$ 検定で $p < 0.001$）。1{,}823 ケースのうち、微分方程式から規則で式を足して作った 1{,}000 ケース（正解式が平均 5.5 本。全ケースでは 3.9 本）だけで評価しても、同じ傾向であった（Table~\ref{tab:result} の右の列）。

説明文で照合する方法が弱い主な原因は、汎用の E5 が似ているが別の量を区別できず（Table~\ref{tab:compare} の例 3）、正解式が第 1 段で候補から落ちることにある。第 1 段の候補（$N_1 = 50$ 件）に正解式が残る割合は、記号で照合すると 0.969、説明文では 0.710 であった（全ケース）。また、利用者側の説明文は正解式の説明を LLM で言い換えて作ったので、利用者が一から書いた説明文では結果がさらに下がる可能性がある。


\section{今後の計画}\label{sec:plan}
\begin{itemize}[leftmargin=1.5em]
  \item \textbf{10 月 16 日まで}：記号が一致するときは記号で、一致しないときは説明文で照合する併用を試す。
  \item \textbf{次回ゼミまでに}：汎用の E5 は似ているが別の量を区別できなかったので、変数の定義が同じ量を指すかを判定する言語モデル ProcessBERT~\citep{ProcessBERT}で説明文を照合する方法を試す。
\end{itemize}


%----------------------------------------------------------------------
% References（bib ファイルを使う。表示順は本文の初出順）
\bibliographystyle{unsrtnat}
\bibliography{progress_report_2026-10-08_var_description}

%=============================%
\end{document}
